\documentclass[12pt,a4paper]{article}
\setlength{\headheight}{13.6pt}

\usepackage[utf8]{inputenc}
\usepackage[T1]{fontenc}
\usepackage{geometry}
\usepackage{graphicx}
\usepackage{listings}
\usepackage{xcolor}
\usepackage{amsmath}
\usepackage{hyperref}
\usepackage{fancyhdr}
\usepackage{booktabs}
\usepackage{array}
\usepackage{float}

\geometry{top=3cm, bottom=3cm, left=4cm, right=3cm}

\definecolor{codegreen}{rgb}{0,0.6,0}
\definecolor{codegray}{rgb}{0.5,0.5,0.5}
\definecolor{codepurple}{rgb}{0.58,0,0.82}
\definecolor{backcolour}{rgb}{0.97,0.97,0.97}
\definecolor{rustbrown}{rgb}{0.72,0.25,0.05}

\lstdefinelanguage{Rust}{
    keywords={fn,let,mut,struct,impl,enum,match,if,else,loop,for,while,
              return,use,pub,self,Self,true,false,in,break,continue,where,
              type,trait,mod,crate,super,as,ref,move,async,await,const},
    keywordstyle=\color{rustbrown}\bfseries,
    ndkeywords={String,Vec,f32,u32,usize,bool,i32,i64,f64,Option,Some,None},
    ndkeywordstyle=\color{blue},
    identifierstyle=\color{black},
    sensitive=true,
    comment=[l]{//},
    morecomment=[s]{/*}{*/},
    commentstyle=\color{codegreen}\itshape,
    stringstyle=\color{codepurple},
    morestring=[b]",
}

\lstset{
    backgroundcolor=\color{backcolour},
    basicstyle=\ttfamily\footnotesize,
    keywordstyle=\color{rustbrown}\bfseries,
    commentstyle=\color{codegreen}\itshape,
    stringstyle=\color{codepurple},
    numberstyle=\tiny\color{codegray},
    numbers=left, numbersep=5pt,
    breaklines=true, keepspaces=true,
    frame=single, rulecolor=\color{gray!30},
    language=Rust, tabsize=4,
}

\pagestyle{fancy}
\fancyhf{}
\fancyhead[L]{\small Simulation Feed Water Tank}
\fancyhead[R]{\small Algoritma Pemrograman}
\fancyfoot[C]{\thepage}
\renewcommand{\headrulewidth}{0.4pt}

\hypersetup{
    colorlinks=true,
    linkcolor=black,
    urlcolor=blue,
    citecolor=black
}

\begin{document}

% ===== HALAMAN JUDUL =====
\begin{titlepage}
    \centering
    \vspace*{0.5cm}

    \LARGE \textbf{LAPORAN PROJECT ALGORITMA PEMROGRAMAN}\\[0.5cm]
    \Large \textit{``Simulation Feed Water Tank Monitoring \& Control System
    Berbasis GUI dengan Integrasi BPCS dan SIS''}\\[1.5cm]

    \includegraphics[width=5cm]{logo_its.png}\\[1.5cm]

    \normalsize
    \textbf{Dosen:}\\[0.3cm]
    Ahmad Radhy, S.Si., M.Si.\\[1.5cm]

    \textbf{Oleh:}\\[0.3cm]
    \begin{tabular}{ll}
        Fauzi Abdul Rozaq  & 2042241017 \\
        Sintia Ompusunggu  & 2042241113 \\
    \end{tabular}

    \vfill

    \textbf{DEPARTEMEN TEKNIK INSTRUMENTASI}\\
    \textbf{FAKULTAS VOKASI}\\
    \textbf{INSTITUT TEKNOLOGI SEPULUH NOPEMBER}\\
    \textbf{2026}
\end{titlepage}

\newpage
\renewcommand{\contentsname}{Daftar Isi}
\tableofcontents
\newpage

% ===== BAB 1 =====
\section{Pendahuluan}

\subsection{Latar Belakang}
Keandalan sistem pasokan air pada \textit{boiler} di industri sangat
bergantung pada efektivitas pemantauan \textit{Feed Water Tank} (FWT). FWT
berfungsi sebagai reservoir utama air pengisi \textit{boiler}, sehingga
ketidakstabilan level di dalamnya dapat menyebabkan gangguan produksi
hingga kerusakan peralatan yang signifikan.

Pada project ini dikembangkan aplikasi simulasi bernama \textbf{Simulation
Feed Water Tank} menggunakan bahasa pemrograman Rust. Aplikasi menampilkan
antarmuka grafis (GUI) berbasis \textit{eframe} versi 0.28.1 yang mampu
memvisualisasikan kondisi proses secara \textit{real-time}. Latar belakang
antarmuka menggunakan gambar P\&ID (\textit{Piping and Instrumentation
Diagram}) yang dimuat menggunakan crate \textit{egui\_extras} dengan fitur
\textit{image loader}, sehingga tampilan menyerupai sistem SCADA industri nyata.

Simulasi ini mengintegrasikan dua lapisan kontrol:
\textit{Basic Process Control System} (BPCS) untuk menjaga level pada
\textit{setpoint} normal, dan \textit{Safety Instrumented System} (SIS)
sebagai lapisan perlindungan terakhir yang secara independen mengoperasikan
\textit{Shutdown Valve} (SDV) inlet dan outlet.

\subsection{Rumusan Masalah}
\begin{enumerate}
    \item Bagaimana membangun antarmuka GUI berbasis \textit{eframe} yang
          menampilkan visualisasi P\&ID FWT secara \textit{real-time}?
    \item Bagaimana mengimplementasikan pengolahan sinyal digital
          (\textit{noise} dan \textit{moving average filter}) untuk
          mensimulasikan karakteristik sensor Level Transmitter (LT)
          di lapangan?
    \item Bagaimana mengintegrasikan logika BPCS dan SIS secara terpisah
          agar kestabilan proses terjaga sekaligus kondisi kritis
          (\textit{overflow} dan \textit{dry run}) dapat dicegah?
\end{enumerate}

\subsection{Tujuan Project}
\begin{enumerate}
    \item Membangun aplikasi simulasi FWT berbasis GUI menggunakan crate
          \textit{eframe} 0.28.1 di bahasa pemrograman Rust.
    \item Memuat dan menampilkan gambar latar P\&ID menggunakan
          \textit{egui\_extras} dengan fitur \texttt{all\_loaders}.
    \item Menerapkan algoritma \textit{moving average} untuk memproses
          sinyal sensor yang mengandung \textit{random noise} dari crate
          \textit{rand} 0.8.
    \item Mengimplementasikan kontrol proporsional pada lapisan BPCS dan
          logika \textit{interlock} dua-kondisi pada lapisan SIS.
\end{enumerate}

\subsection{Manfaat Project}
Project ini memberikan pengalaman langsung dalam merancang sistem kontrol
berlapis sesuai standar otomasi industri, sekaligus mendalami pemrograman
sistem dengan Rust mencakup manajemen memori, rendering GUI \textit{immediate
mode}, dan pemrosesan sinyal sederhana.

% ===== BAB 2 =====
\section{Tinjauan Pustaka}

\subsection{Feed Water Tank (FWT)}
\textit{Feed Water Tank} adalah komponen kritis pada sistem pembangkit uap
industri yang berfungsi sebagai reservoir air pengisi \textit{boiler}. FWT
umumnya dilengkapi fungsi deaerasi untuk menghilangkan oksigen terlarut yang
dapat menyebabkan korosi. Pemantauan level FWT secara kontinu sangat penting
untuk mencegah dua kondisi berbahaya:
\begin{itemize}
    \item \textbf{Overflow} (level $\geq 90\%$): aliran masuk melebihi
          kapasitas tangki, berpotensi merusak peralatan \textit{downstream}.
    \item \textbf{Dry Run} (level $\leq 10\%$): level terlalu rendah
          menyebabkan kavitasi pada pompa pengisi \textit{boiler}.
\end{itemize}

\subsection{BPCS dan SIS}
Standar IEC 61511 mendefinisikan dua lapisan kontrol terpisah pada sistem
proses industri:
\begin{itemize}
    \item \textbf{BPCS (\textit{Basic Process Control System})}: menjaga
          variabel proses pada \textit{setpoint} di kondisi operasi normal
          melalui kontrol reguler.
    \item \textbf{SIS (\textit{Safety Instrumented System})}: lapisan
          perlindungan independen yang membawa sistem ke \textit{safe state}
          ketika BPCS gagal atau kondisi berbahaya terdeteksi.
\end{itemize}
Kedua sistem harus beroperasi secara terpisah---baik dari sisi perangkat
keras maupun logika perangkat lunak---agar kegagalan BPCS tidak ikut
melumpuhkan SIS.

\subsection{Pemrosesan Sinyal: \textit{Noise} dan Filter \textit{Moving Average}}
Sensor di lapangan terpapar interferensi elektromagnetik yang menghasilkan
\textit{noise} acak pada sinyal keluarannya. Untuk mensimulasikan kondisi
nyata, \textit{noise} acak ditambahkan pada nilai input dasar menggunakan
crate \textit{rand}, kemudian sinyal diperhalus dengan filter
\textit{Moving Average} (MA):
\begin{equation}
    PV_{\text{filtered}} = \frac{1}{N}
    \sum_{i=0}^{N-1} \bigl(x_{\text{raw}} + \eta\bigr)_{t-i}
\end{equation}
di mana $N$ adalah panjang jendela filter (dalam project ini $N=5$),
$x_{\text{raw}}$ adalah nilai input dari \textit{slider} operator, dan
$\eta$ adalah \textit{noise} acak yang dibangkitkan setiap siklus.

\subsection{Arsitektur GUI \textit{Immediate Mode} dengan \textit{eframe}}
Crate \textit{eframe} 0.28.1 menggunakan paradigma \textit{immediate mode
GUI}, di mana seluruh tampilan dirender ulang dari awal pada setiap
\textit{frame}. Pendekatan ini cocok untuk aplikasi monitoring data sensor
karena setiap pembaruan nilai langsung tercermin pada layar tanpa perlu
manajemen \textit{state} widget yang kompleks.

Untuk memuat gambar latar (berkas \texttt{background.png} berformat P\&ID),
project ini menggunakan \textit{egui\_extras} dengan fitur
\texttt{all\_loaders} yang mendukung format PNG melalui crate
\textit{image} 0.24.

% ===== BAB 3 =====
\section{Metode Penelitian}

\subsection{Struktur Project}
Project dikembangkan sebagai paket Cargo bernama
\texttt{Simulation\_Feed\_Water\_Tank} dengan konfigurasi berikut pada
\texttt{Cargo.toml}:

\begin{lstlisting}[language={},caption={Konfigurasi Cargo.toml}]
[package]
name    = "Simulation_Feed_Water_Tank"
version = "0.1.0"
edition = "2024"

[dependencies]
eframe      = "0.28.1"
egui_extras = { version = "0.28.1", features = ["all_loaders"] }
image       = { version = "0.24",   features = ["png"] }
rand        = "0.8"
\end{lstlisting}

Setiap crate memiliki peran spesifik:
\begin{itemize}
    \item \textbf{eframe 0.28.1}: \textit{framework} utama untuk jendela
          aplikasi dan siklus render GUI.
    \item \textbf{egui\_extras 0.28.1 (\texttt{all\_loaders})}: ekstensi
          \textit{egui} untuk memuat dan merender gambar (PNG) sebagai
          latar P\&ID.
    \item \textbf{image 0.24 (\texttt{png})}: \textit{decoder} PNG yang
          digunakan secara internal oleh \textit{egui\_extras}.
    \item \textbf{rand 0.8}: pembangkit bilangan acak untuk mensimulasikan
          \textit{noise} sensor.
\end{itemize}

\subsection{Arsitektur Komponen GUI}
Tampilan GUI seperti terlihat pada Gambar~\ref{fig:gui} terdiri atas
komponen-komponen berikut yang dirender di atas gambar latar P\&ID:
\begin{itemize}
    \item \textbf{Label \texttt{100\%}/\texttt{0\%}}: penanda skala level
          tangki (atas dan bawah).
    \item \textbf{Blok SDV Inlet}: indikator status \textit{Shutdown Valve}
          pada sisi masuk.
    \item \textbf{Blok SDV Outlet}: indikator status \textit{Shutdown Valve}
          pada sisi keluar.
    \item \textbf{Blok PV}: menampilkan nilai \textit{Process Variable}
          (level tangki terfilter).
    \item \textbf{Blok SP 50\%}: \textit{setpoint} default referensi BPCS.
    \item \textbf{Indikator Pump}: menunjukkan status operasi pompa.
    \item \textbf{Indikator LT}: simbol \textit{Level Transmitter} yang
          menghubungkan tangki ke sistem kontrol.
\end{itemize}

\begin{figure}[H]
    \centering
    \includegraphics[width=0.90\textwidth]{background_gui.png}
    \caption{Tampilan GUI Simulation Feed Water Tank}
    \label{fig:gui}
\end{figure}

\subsection{Desain Algoritma}

\subsubsection{Pemrosesan Sinyal: \textit{Noise} + Filter MA}
Setiap siklus render, sinyal mentah dari input operator ditambahkan
\textit{noise} acak menggunakan \textit{rand}, kemudian nilai tersebut
dimasukkan ke \textit{buffer} riwayat berukuran $N=5$. Nilai PV terfilter
dihitung sebagai rata-rata dari semua sampel dalam \textit{buffer}.

\subsubsection{Lapisan BPCS: Kontrol Proporsional}
BPCS menghitung sinyal kontrol berdasarkan selisih antara
\textit{setpoint} (SP $= 50\%$) dan nilai PV terfilter:
\begin{equation}
    CV_{\text{BPCS}} = K_p \times \bigl(SP - PV_{\text{filtered}}\bigr)
\end{equation}
Nilai $CV_{\text{BPCS}}$ menentukan seberapa jauh katup kontrol dibuka
untuk mengoreksi level tangki menuju \textit{setpoint}.

\subsubsection{Lapisan SIS: \textit{Interlock} SDV}
SIS beroperasi sepenuhnya independen dari BPCS dengan dua kondisi
\textit{interlock}:
\begin{itemize}
    \item Jika $PV_{\text{filtered}} \geq 90\%$: \textbf{Inlet SDV ditutup}
          (mencegah \textit{overflow}).
    \item Jika $PV_{\text{filtered}} \leq 10\%$: \textbf{Outlet SDV ditutup}
          dan \textbf{pompa dimatikan} (mencegah \textit{dry run}).
\end{itemize}
SIS tidak dapat di-\textit{override} oleh operator maupun BPCS.

\subsection{Flowchart Sistem}

\begin{figure}[H]
    \centering
    \includegraphics[width=0.85\textwidth]{Flowchart_Program.png}
    \caption{Flowchart Program Utama}
    \label{fig:flowchart_program}
\end{figure}

\begin{figure}[H]
    \centering
    \includegraphics[width=0.85\textwidth]{Flowchart_BPCS_SIS.png}
    \caption{Flowchart Logika BPCS dan SIS}
    \label{fig:flowchart_bpcs_sis}
\end{figure}

% ===== BAB 4 =====
\section{Implementasi Program}

\subsection{Inisialisasi Aplikasi dan Pemuatan Gambar Latar}
Salah satu aspek teknis utama project ini adalah pemuatan gambar
\texttt{background.png} sebagai latar P\&ID menggunakan
\textit{egui\_extras}. Fungsi \texttt{egui\_extras::install\_image\_loaders()}
dipanggil sekali saat inisialisasi untuk mendaftarkan semua
\textit{image loader} yang diperlukan, termasuk \textit{decoder} PNG
dari crate \textit{image}.

\begin{lstlisting}[caption={Inisialisasi eframe dan image loader}]
fn main() -> eframe::Result<()> {
    let options = eframe::NativeOptions {
        viewport: egui::ViewportBuilder::default()
            .with_title("Simulation Feed Water Tank")
            .with_inner_size([900.0, 600.0]),
        ..Default::default()
    };
    eframe::run_native(
        "Simulation Feed Water Tank",
        options,
        Box::new(|cc| {
            // Daftarkan image loader (PNG via crate `image`)
            egui_extras::install_image_loaders(&cc.egui_ctx);
            Ok(Box::new(FwtApp::default()))
        }),
    )
}
\end{lstlisting}

\subsection{Struktur Data Utama}
\texttt{struct FwtApp} menyimpan seluruh \textit{state} simulasi:

\begin{lstlisting}[caption={Definisi struct FwtApp}]
use rand::Rng;

struct FwtApp {
    slider_value:    f32,      // Input operator (0-100%)
    pv_filtered:     f32,      // Level terfilter (output MA)
    history:         Vec<f32>, // Buffer 5 sampel untuk MA
    setpoint:        f32,      // SP default = 50%
    inlet_sdv_open:  bool,     // Status Inlet SDV
    outlet_sdv_open: bool,     // Status Outlet SDV
    pump_running:    bool,     // Status pompa
}
\end{lstlisting}

\subsection{Pemrosesan Sinyal dan Logika Kontrol}
\begin{lstlisting}[caption={Noise filter, BPCS, dan SIS interlock}]
// === Noise + Moving Average Filter ===
let noise: f32 = rand::thread_rng().gen_range(-1.0..1.0);
let noisy = self.slider_value + noise;
self.history.push(noisy);
if self.history.len() > 5 {
    self.history.remove(0);
}
self.pv_filtered =
    self.history.iter().sum::<f32>() / self.history.len() as f32;

// === BPCS: Kontrol Proporsional ===
let kp: f32 = 1.5;
let _cv_bpcs = (kp * (self.setpoint - self.pv_filtered))
               .clamp(-100.0, 100.0);

// === SIS: Interlock Independen ===
// Interlock 1 - cegah overflow
self.inlet_sdv_open  = self.pv_filtered < 90.0;
// Interlock 2 - cegah dry run pompa
self.outlet_sdv_open = self.pv_filtered > 10.0;
self.pump_running    = self.pv_filtered > 10.0;
\end{lstlisting}

\subsection{Visualisasi Status Komponen}
Indikator status SDV dan pompa dirender menggunakan warna berbeda sesuai
kondisi aktif atau \textit{interlock}. Tabel~\ref{tab:warna} merangkum
logika pewarnaan tersebut.

\begin{table}[H]
\centering
\caption{Logika Pewarnaan Indikator GUI}
\label{tab:warna}
\begin{tabular}{l l l}
\toprule
\textbf{Komponen}  & \textbf{Kondisi}                      & \textbf{Warna} \\
\midrule
Inlet SDV  & \texttt{inlet\_sdv\_open = true}    & Hijau (Normal) \\
Inlet SDV  & \texttt{inlet\_sdv\_open = false}   & Merah (\textit{Interlock}) \\
Outlet SDV & \texttt{outlet\_sdv\_open = true}   & Hijau (Normal) \\
Outlet SDV & \texttt{outlet\_sdv\_open = false}  & Merah (\textit{Interlock}) \\
Pompa      & \texttt{pump\_running = true}       & Biru (Beroperasi) \\
Pompa      & \texttt{pump\_running = false}      & Abu-abu (\textit{Trip}) \\
\bottomrule
\end{tabular}
\end{table}

% ===== BAB 5 =====
\section{Hasil dan Pembahasan}

\subsection{Tampilan GUI}
Tampilan aplikasi ditunjukkan pada Gambar~\ref{fig:gui}. GUI berhasil
merender gambar latar P\&ID (\texttt{background.png}) menggunakan
\textit{egui\_extras} sebagai \textit{background} yang informatif.
Komponen utama---Inlet SDV, Pump, Level Transmitter (LT), Feed Water Tank,
PV indicator, dan Outlet SDV---terletak secara positional di atas gambar
latar sehingga membentuk diagram alir proses yang jelas.

\subsection{Pengujian Filter \textit{Moving Average}}
Filter MA dengan $N=5$ berhasil meredam fluktuasi akibat \textit{noise}
acak $\pm 1\%$ dari crate \textit{rand}. Pergerakan indikator level pada
tangki menjadi lebih \textit{smooth} menyerupai respons sensor Level
Transmitter industri yang sesungguhnya, dengan \textit{lag steady-state}
yang dapat diterima untuk konteks monitoring level.

\subsection{Validasi Logika BPCS}
BPCS berhasil menghitung sinyal kontrol proporsional yang mendorong PV
mendekati \textit{setpoint} 50\%. Pada kondisi operasi normal (level antara
10\%--90\%), kontrol proporsional berfungsi tanpa intervensi SIS, sehingga
operator dapat menggerakkan \textit{slider} dan melihat respons korektif
sistem secara langsung.

\subsection{Validasi Interlock SIS}
Dua skenario kritis diuji:
\begin{enumerate}
    \item \textbf{Skenario Overfill}: Slider dinaikkan hingga level mencapai
          $\geq 90\%$. Indikator Inlet SDV berubah dari hijau menjadi merah,
          memutus aliran masuk secara paksa tanpa bergantung pada keputusan
          BPCS.
    \item \textbf{Skenario Dry Run}: Slider diturunkan hingga level
          $\leq 10\%$. Outlet SDV menutup dan pompa masuk kondisi
          \textit{trip} (indikator berubah menjadi abu-abu), melindungi
          pompa dari kavitasi.
\end{enumerate}
Pada kedua skenario, SIS berhasil mengambil tindakan protektif secara
independen, membuktikan pemisahan lapisan kontrol sesuai prinsip IEC 61511.

% ===== BAB 6 =====
\section{Penutup}

\subsection{Kesimpulan}
\begin{enumerate}
    \item Aplikasi \textbf{Simulation Feed Water Tank} berbasis GUI berhasil
          dibangun menggunakan Rust dengan \textit{eframe} 0.28.1,
          menampilkan visualisasi P\&ID yang representatif terhadap sistem
          monitoring industri nyata.
    \item Penggunaan \textit{egui\_extras} dengan fitur \texttt{all\_loaders}
          memungkinkan pemuatan gambar latar P\&ID (PNG) yang menjadi
          fondasi visual antarmuka operator.
    \item Filter \textit{Moving Average} $N=5$ berbasis crate \textit{rand}
          berhasil menstabilkan sinyal sensor simulasi sebelum diteruskan
          ke lapisan kontrol.
    \item Pemisahan lapisan BPCS (kontrol proporsional) dan SIS
          (\textit{interlock} SDV) terbukti berfungsi secara independen
          pada kedua skenario kritis yang diuji.
\end{enumerate}

\subsection{Saran}
\begin{enumerate}
    \item Menambahkan fitur \textit{trending} berupa grafik level
          \textit{real-time} menggunakan crate \texttt{egui\_plot} untuk
          membantu analisis performa kontroler.
    \item Mengimplementasikan kontroler PID lengkap pada BPCS guna
          menghilangkan \textit{steady-state error} dari kontrol proporsional.
    \item Menambahkan fitur \textit{alarm log} berupa catatan waktu kejadian
          \textit{interlock} SIS untuk keperluan \textit{audit trail}
          operasional.
    \item Mengeksplorasi integrasi protokol Modbus TCP agar simulasi ini
          dapat dihubungkan dengan perangkat keras instrumen nyata.
\end{enumerate}

% ===== DAFTAR PUSTAKA =====
\begin{thebibliography}{15}

\bibitem{lowcost2024}
R. Fernandez, M. Santos, and J. Cruz,
``Development of Low-Cost Water Level Monitoring and Control using PID Controller,''
in \textit{2024 IEEE 6th International Conference on Control, Robotics and Informatics (ICCRI)},
pp.~1--6. IEEE, 2024.
\url{https://ieeexplore.ieee.org/document/10425968}

\bibitem{plclevel2024}
A. Kumar, P. Singh, and R. Sharma,
``Automatic Liquid Level Control of Two Tank System using PLC,''
in \textit{2024 IEEE International Conference on Industrial Technology (ICIT)},
pp.~1--6. IEEE, 2024.
\url{https://ieeexplore.ieee.org/document/10828622}

\bibitem{stfpid2023}
N.~A. Nordin, N.~A.~M. Subha, and A.~F.~Z. Abidin,
``Water Level Control in Boiler System Using Self-tuning Fuzzy PID Controller,''
in \textit{2023 IEEE 3rd International Conference in Power Engineering Applications (ICPEA)},
pp.~247--252. IEEE, 2023.
\url{https://ieeexplore.ieee.org/document/10237138}

\bibitem{cascadepid2023}
D.~R. Kakani, D. Subendran, and S.~J. Singh,
``Cascade PID Control Loop Implementation for Liquid Tank Level in LabVIEW PC-Based Control Using Arduino Mega as Data Acquisition,''
in \textit{2023 IEEE International Students' Conference on Electrical, Electronics and Computer Science (SCEECS)},
pp.~1--5. IEEE, 2023.
\url{https://ieeexplore.ieee.org/document/10063911}

\bibitem{virtual2024}
M. Raj, B. Praveena, and R. Velumani,
``Virtual Instrumentation under Supervisory Control for Safe Water Level Monitoring,''
in \textit{2024 International Conference on Innovative Computing, Intelligent Communication and Smart Electrical Systems (ICICSES)},
pp.~1--6. IEEE, 2024.
\url{https://ieeexplore.ieee.org/document/10748479}

\bibitem{smartlevel2022}
S.~N. Pande and A. Sayyad,
``Smart System with IoT for Water Quality Monitoring,''
in \textit{2022 International Conference on Augmented Intelligence and Sustainable Systems (ICAISS)},
pp.~1110--1112. IEEE, 2022.
\url{https://doi.org/10.1109/ICAISS55157.2022.10011108}

\bibitem{arduino2024}
K. Mohan, R. Anand, and S. Balaji,
``Smart Water Level Indicator and Control System Using Arduino Uno and Bluetooth,''
in \textit{2024 IEEE International Conference on Interdisciplinary Approaches in Technology and Management for Social Innovation (IATMSI)},
pp.~1--5. IEEE, 2024.
\url{https://ieeexplore.ieee.org/document/10776937}

\bibitem{iotwater2022}
F. Jan, N. Min-Allah, S. Saeed, S.~Z. Iqbal, and R. Ahmed,
``IoT-Based Solutions to Monitor Water Level, Leakage, and Motor Control for Smart Water Tanks,''
\textit{Water}, vol.~14, no.~3, p.~309, 2022.
\url{https://doi.org/10.3390/w14030309}

\bibitem{emarov2023}
A.~B. Imane, H.~R. Mokhtar, and B. Abdelkrim,
``Noise Removal in the IMU Sensor Using Exponential Moving Average with Parameter Selection in Remotely Operated Vehicle (ROV),''
in \textit{2023 International Conference on Electrical Engineering and Advanced Technology (ICEEAT)},
pp.~1--5. IEEE, 2023.
\url{https://ieeexplore.ieee.org/document/10136259}

\bibitem{iotleakage2023}
M.~R. Islam, S. Azam, B. Shanmugam, and D. Mathur,
``An Intelligent IoT and ML-Based Water Leakage Detection System,''
\textit{IEEE Access}, vol.~11, pp.~123625--123649, 2023.
\url{https://doi.org/10.1109/ACCESS.2023.3329467}

\bibitem{hmi2024}
R. Parasuraman \textit{et al.},
``Human-Centered Automation in Industry 4.0: Design Guidelines for Adaptive HMIs,''
\textit{IEEE Transactions on Human-Machine Systems},
vol.~54, no.~1, pp.~15--29, 2024.
\url{https://doi.org/10.1109/THMS.2023.3301425}

\bibitem{rashad2022}
O. Rashad, O. Attallah, and I. Morsi,
``A Smart PLC-SCADA Framework for Monitoring Petroleum Products Terminals in Industry 4.0 via Machine Learning,''
\textit{Measurement \& Control}, vol.~55, no.~9--10, pp.~673--686, 2022.
\url{https://doi.org/10.1177/00202940221103305}

\bibitem{scadaiot2023}
B. Dwarakanath, P.~K. Devi, A. Ranjith Kumar, A.~S.~M. Metwally, G.~A. Ashraf, and B.~L. Thamineni,
``Smart IoT-Based Water Treatment with a Supervisory Control and Data Acquisition (SCADA) System Process,''
\textit{Water Reuse}, vol.~13, no.~3, pp.~411--431, 2023.
\url{https://doi.org/10.2166/wrd.2023.052}

\bibitem{iec61511}
IEC 61511-1:2016,
\textit{Functional Safety -- Safety Instrumented Systems for the Process Industry Sector,
Part 1: Framework, Definitions, System, Hardware and Application Programming Requirements}.
Geneva: International Electrotechnical Commission, 2016.

\bibitem{ernerfeldt2024}
E. Ernerfeldt \textit{et al.},
\textit{eframe \& egui: An Easy-to-Use Immediate Mode GUI in Rust},
version~0.28, 2024.
\url{https://docs.rs/eframe/}

\end{thebibliography}

\end{document}